\section{Circuit Complexity --- Uniform Tableau Emitter Main}
%%% generated by the blueprint pipeline from TCSlib.Complexity.CircuitComplexity.UniformTableauEmitterMain
\subsection*{Circuit Complexity --- Uniform Tableau Emitter Main: Introduction}

Declarations of the Lean module \texttt{TCSlib.Complexity.CircuitComplexity.UniformTableauEmitterMain}.

\subsection*{Circuit Complexity --- Uniform Tableau Emitter Main: Results}

\begin{theorem}[Unary input scan]
\lean{Complexity.UTab.goes_rN}
\label{Complexity.UTab.goes_rN}
\leanok
\uses{Complexity.CounterProg.Goes, Complexity.CounterProg.Goes.congr, Complexity.CounterProg.Goes.trans, Complexity.CounterProg.goes_out, Complexity.CounterProg.goes_rd_true, Complexity.UTab.ER, Complexity.UTab.ER.f, Complexity.UTab.goes_incR, Complexity.UTab.goes_refl, Complexity.UTab.prog, Complexity.UTab.rNC}
\difficulty{6}
Let $M$ be a finite-state Turing machine with Boolean input, let $w$, $pre$, and $rest$
be Boolean lists, let $a$ be a natural number, and let $e$ be an eleven-register
configuration record. If $w = pre \, ++ \, (1^a \, ++ \, rest)$, then for every prior
output list $o$, the counter program for $M$, started at label $rN$, with registers
given by $e$, input position $|pre|$, and output $o$, reaches label $rN$ with the
register record obtained from $e$ by increasing its input-count register by $a$, input
position $|pre|+a$, and output $o \, ++ \, 1^a$, in at most $3a$ steps.
\end{theorem}

\begin{theorem}[Unary step-count scan]
\lean{Complexity.UTab.goes_rT}
\label{Complexity.UTab.goes_rT}
\leanok
\uses{Complexity.CounterProg.Goes, Complexity.CounterProg.Goes.congr, Complexity.CounterProg.Goes.trans, Complexity.CounterProg.goes_rd_true, Complexity.UTab.ER, Complexity.UTab.ER.f, Complexity.UTab.goes_incR, Complexity.UTab.goes_refl, Complexity.UTab.prog, Complexity.UTab.rT}
\difficulty{5}
Let $M$ be a finite-state Turing machine with Boolean input, let $w$, $pre$, and $rest$
be Boolean lists, let $b$ be a natural number, and let $e$ be an eleven-register
configuration record. If $w = pre \, ++ \, (1^b \, ++ \, rest)$, then for every prior
output list $o$, the counter program for $M$, started at label $rT$, with registers
given by $e$, input position $|pre|$, and output $o$, reaches label $rT$ with the
register record obtained from $e$ by increasing its step-count register by $b$, input
position $|pre|+b$, and unchanged output $o$, in at most $2b$ steps.
\end{theorem}

\begin{theorem}[Hard-wired input scan]
\lean{Complexity.UTab.goes_rZ}
\label{Complexity.UTab.goes_rZ}
\leanok
\uses{BoolCircuit.constGate, Complexity.CounterProg.Goes, Complexity.CounterProg.Goes.congr, Complexity.CounterProg.Goes.trans, Complexity.CounterProg.flatMap_exec_bitsOps, Complexity.CounterProg.goes_rd_false, Complexity.CounterProg.goes_rd_true, Complexity.UTab.ER, Complexity.UTab.ER.f, Complexity.UTab.Lmax, Complexity.UTab.gbits, Complexity.UTab.goes_incR, Complexity.UTab.goes_refl, Complexity.UTab.goes_tp, Complexity.UTab.length_tmpl_le, Complexity.UTab.prog, Complexity.UTab.rZ, Complexity.UTab.tmpl, Complexity.UTab.tpL}
\difficulty{6}
Let $M$ be a finite-state Turing machine with Boolean input, let $w$, $pre$, and $rest$
be Boolean lists, let $z$ be a Boolean list, and let $e$ be an eleven-register
configuration record. If $w = pre \, ++ \, (z \, ++ \, rest)$, then for every prior
output list $o$, the counter program for $M$, started at label $rZ$, with registers
given by $e$, input position $|pre|$, and output $o$, reaches label $rZ$ with the
register record obtained from $e$ by increasing its hard-wired-bit count by $|z|$, input
position $|pre|+|z|$, and output $o$ followed by the concatenation, for each bit $b$ of
$z$ in order, of the code consisting of a true bit followed by the encoding of the
constant gate for $b$. The run takes at most $|z|(L_{\max}(M)+3)$ steps.
\end{theorem}

\begin{theorem}[The stop instruction halts immediately]
\lean{Complexity.UTab.goes_stop}
\label{Complexity.UTab.goes_stop}
\leanok
\uses{Complexity.CounterProg.Goes, Complexity.CounterProg.goes_halt, Complexity.UTab.prog}
\difficulty{1}
Let $M$ be a finite-state Turing machine with Boolean symbols, let $w$ be a finite
Boolean input, let $\rho$ assign a natural-number value to each of eleven registers, and
let $p$ be a natural-number input position. Starting the associated counter program at
its stop label with register assignment $\rho$, input position $p$, and any previously
printed output, it halts in one step, leaving the registers and input position unchanged
and appending no output.
\end{theorem}

\begin{theorem}[Range expansion of program gates]
\lean{Complexity.UTab.tabProgGates_range'}
\label{Complexity.UTab.tabProgGates_range'}
\leanok
\uses{BoolCircuit.GInstr, BoolCircuit.tabInstrGates, BoolCircuit.tabProgGates, Complexity.UTab.range'_succ_left}
\difficulty{4}
Let $m,w,n,L,a\in\mathbb{N}$, let $F$ assign a bit string to each element of a finite
type, let $z$ be a bit string, and let $f$ assign a gate instruction to each natural
number. The gates produced by the tabulated program starting at $a$ from the list of
instructions $f(a),\ldots,f(a+L-1)$ are the concatenation, in increasing index order, of
the gates produced separately by each instruction at its corresponding index.
\end{theorem}

\begin{theorem}[Configuration tableau output index]
\lean{Complexity.UTab.cfgTab_output}
\label{Complexity.UTab.cfgTab_output}
\leanok
\uses{BoolCircuit.tabBase, BoolCircuit.tabCircuit, Complexity.CounterProg.LinE.val, Complexity.UTab.ER.f, Complexity.UTab.ER.f_i, Complexity.UTab.ER.f_nc, Complexity.UTab.ER.f_z, Complexity.UTab.SSrc.toLinE, Complexity.UTab.Str, Complexity.UTab.lstart, Complexity.cfgProg, Complexity.cfgStride, Complexity.cfgTab, Complexity.cfgWidth, Complexity.length_cfgProg, Complexity.length_tabLayout, Complexity.snapIdx, Complexity.snapIdx_lt_succ, Complexity.snapWidth, Complexity.tabLayout}
\difficulty{3}
Let $M$ be a finite-state Turing machine with Boolean work-tape symbols, let
$n,T\in\mathbb{N}$, and let $z$ be a finite Boolean string. The output vertex of the
configuration-tableau circuit for $M$, with $n$ circuit inputs, hard-wired string $z$,
and time bound $T$, is the value of the linear expression for the current snapshot’s
encoded output bit, evaluated on the register values at time $T+1$ with loop counter
zero.
\end{theorem}

\begin{theorem}[Configuration tableau encoding]
\lean{Complexity.UTab.cfgTab_encode}
\label{Complexity.UTab.cfgTab_encode}
\leanok
\uses{BoolCircuit.DAGCircuit.encode, BoolCircuit.constGate, BoolCircuit.encodeNat, BoolCircuit.tabCircuit, BoolCircuit.tabPre, Complexity.CounterProg.LinE.val, Complexity.UTab.ER.f, Complexity.UTab.SSrc.toLinE, Complexity.UTab.Str, Complexity.UTab.Wd, Complexity.UTab.cfgTab_output, Complexity.UTab.encodeList_eq_gbits, Complexity.UTab.gbits, Complexity.UTab.gbits_append, Complexity.UTab.ibits, Complexity.UTab.lstart, Complexity.UTab.tabProgGates_range', Complexity.cfgProg, Complexity.cfgStride, Complexity.cfgTab, Complexity.snapWidth, Complexity.tabLayout}
\difficulty{3}
Let $M$ be a finite-state Turing machine with Boolean work-tape symbols, let
$n,T\in\mathbb{N}$, and let $z$ be a finite Boolean string. The Boolean encoding of the
configuration-tableau circuit for $M$, with $n$ circuit inputs, hard-wired string $z$,
and time bound $T$, consists of $n$ initial ones followed by zero, then the encoded
preamble (false and true constant gates, the gates for the bits of $z$, and the
prescribed false constant gates), the encodings of all tableau-program instructions in
order, and a final zero followed by a unary string. The length of that final unary
string is the value of the linear expression for the current snapshot’s encoded output
bit, evaluated on the register values at time $T+1$ with loop counter zero.
\end{theorem}

\begin{definition}[Simulation prefix encoding]
\lean{Complexity.UTab.prefixBits}
\label{Complexity.UTab.prefixBits}
\leanok
\uses{BoolCircuit.constGate, Complexity.UTab.Wd, Complexity.UTab.gbits}
For a finite binary Turing machine $M$, a natural number $n$, and a bit string $z$,
define the prefix bit string by concatenating $n$ copies of $1$, then $0$, then the gate
encoding of the two constant gates for false and true, followed by the gate encodings of
one constant gate for each bit of $z$, and finally the gate encodings of $W_d(M)$
constant-false gates.
\end{definition}

\begin{definition}[Prefix length bound]
\lean{Complexity.UTab.prefixBound}
\label{Complexity.UTab.prefixBound}
\leanok
\uses{Complexity.UTab.Lmax}
Let $M$ be a finite-tape Turing machine with associated quantities $L(M)$, and let $n,T$
be nonnegative integers and $z$ a finite Boolean word. Define the prefix bound to be
$3n+2T+|z|(L(M)+3)+2(L(M)+1)+6$.
\end{definition}

\begin{theorem}[Counter program prefix execution]
\lean{Complexity.UTab.goes_prefix}
\label{Complexity.UTab.goes_prefix}
\leanok
\uses{BoolCircuit.constGate, Complexity.CounterProg.Goes, Complexity.CounterProg.Goes.congr, Complexity.CounterProg.Goes.trans, Complexity.CounterProg.flatMap_exec_bitsOps, Complexity.CounterProg.goes_out, Complexity.CounterProg.goes_rd_end, Complexity.CounterProg.goes_rd_false, Complexity.UTab.ER, Complexity.UTab.ER.f, Complexity.UTab.ER.zero, Complexity.UTab.Lmax, Complexity.UTab.gbits, Complexity.UTab.goes_rN, Complexity.UTab.goes_rT, Complexity.UTab.goes_rZ, Complexity.UTab.goes_tp, Complexity.UTab.length_tmpl_le, Complexity.UTab.lstart, Complexity.UTab.prefixBits, Complexity.UTab.prefixBound, Complexity.UTab.prog, Complexity.UTab.tmpl}
\difficulty{5}
Let $M$ be a Boolean multi-tape Turing machine, let $n,T$ be natural numbers, and let
$z$ be a finite Boolean word. On input consisting of $n$ ones, a zero, $T$ ones, a
second zero, and then $z$, the eleven-register counter program for $M$ has a run of at
most $3n+2T+|z|(L+3)+2(L+1)+6$ steps from its initial label, with all registers and the
input position zero, to the label that begins the layout phase. At that point the input
position is $n+T+|z|+2$ and the registers hold $n$ in the input-count register, $|z|$ in
the hard-wired-length register, and $T$ in the step-budget register, with the other
registers zero. The run appends $n$ ones and a zero, then the encodings of the false and
true constant gates, the constant gates for the successive bits of $z$, and $W$ false
constant gates. Here $W$ is the width of the universal table for the configuration
update rule of $M$, $L$ is the maximum template length for that machine, and each gate
encoding consists of its two-bit gate-kind code followed by a length-prefixed encoding
of its input list; the bitstring encoding of a list places a one before each entry and
ends with a zero.
\end{theorem}

\begin{definition}[Run length bound]
\lean{Complexity.UTab.runBound}
\label{Complexity.UTab.runBound}
\leanok
\uses{Complexity.UTab.Lmax, Complexity.UTab.layerBound, Complexity.UTab.prefixBound}
Let $M$ be a finite-tape Turing machine with associated quantities $L(M)$, and let $n,T$
be nonnegative integers and $z$ a finite Boolean word. Define the run bound as the
prefix bound plus $(T+1)(B+7T+10)+L(M)+2$, where $B$ is the layer bound for $M,n,z,T$.
\end{definition}

\begin{theorem}[Universal Program Emits the Tableau Circuit]
\lean{Complexity.UTab.goes_emit}
\label{Complexity.UTab.goes_emit}
\leanok
\uses{BoolCircuit.DAGCircuit.encode, Complexity.CounterProg.Goes, Complexity.CounterProg.Goes.congr, Complexity.CounterProg.Goes.trans, Complexity.CounterProg.MOp.exec, Complexity.CounterProg.flatMap_exec_linOps, Complexity.UTab.ER.f, Complexity.UTab.ER.zero, Complexity.UTab.cfgTab_encode, Complexity.UTab.goes_layers, Complexity.UTab.goes_prefix, Complexity.UTab.goes_stop, Complexity.UTab.goes_tp, Complexity.UTab.length_tmpl_le, Complexity.UTab.lstart, Complexity.UTab.prefixBits, Complexity.UTab.prog, Complexity.UTab.runBound, Complexity.UTab.tmpl, Complexity.cfgTab}
\difficulty{3}
Let $M$ be a finite-state multitape Turing machine with Boolean input and output, let
$n,T\in\mathbb N$, and let $z$ be a finite bit string, with $T\ge 1$. The counter
program constructed from $M$, started at its initial label with all 11 registers zero
and read position zero, on the input consisting of $n$ ones, then a zero, then $T$ ones,
then a zero, followed by $z$, reaches a halted state with some register assignment and
input position, and has printed the encoding of the tableau circuit for $M$, $n$, $z$,
and $T$; this is achieved within the stated run bound.
\end{theorem}

\begin{definition}[Emission constant]
\lean{Complexity.UTab.emitC}
\label{Complexity.UTab.emitC}
\leanok
\uses{Complexity.UTab.Lmax}
For a finite-tape Turing machine $M$, let $k(M)$ denote its number of tapes and let
$L_{\max}(M)$ denote its associated maximum length. Define the emission constant by
$c(M)=60(k(M)+1)(L_{\max}(M)+30)$.
\end{definition}

\begin{theorem}[Cubic bound on the run length]
\lean{Complexity.UTab.runBound_le}
\label{Complexity.UTab.runBound_le}
\leanok
\uses{Complexity.UTab.Ki, Complexity.UTab.Lmax, Complexity.UTab.Str, Complexity.UTab.emitC, Complexity.UTab.layerBound, Complexity.UTab.prefixBound, Complexity.UTab.runBound, Complexity.cfgStride}
\difficulty{5}
Let $M$ be a finite Turing machine with $k$ work tapes, and let $n,T\in\mathbb{N}$ and
$z$ be a finite binary string. Put $L$ equal to the maximum length of the templates
associated with $M$, $K=L+3$, and $S=k(2T+1)+|z|+n+3$. Define
$A=k((2T+3)(L+20))+(|z|+n+2)(L+20)+K+9(T+1)S+6$ and $B=3n+2T+|z|(L+3)+2(L+1)+6$. Then
$B+(T+1)(A+7T+10)+L+2\le 60(k+1)(L+30)(n+T+|z|+3)^3$.
\end{theorem}

\begin{theorem}[Unary string classification]
\lean{Complexity.UTab.classify}
\label{Complexity.UTab.classify}
\leanok
\difficulty{4}
Every finite Boolean string is either a string consisting entirely of ones, or a string
of zero or more ones followed by a zero and an arbitrary finite Boolean string.
\end{theorem}

\begin{theorem}[Bounded Simulation of a Boolean Turing Machine]
\lean{Complexity.UTab.goes_any}
\label{Complexity.UTab.goes_any}
\leanok
\uses{Complexity.CounterProg.Goes, Complexity.CounterProg.Goes.congr, Complexity.CounterProg.Goes.trans, Complexity.CounterProg.goes_jz_zero, Complexity.CounterProg.goes_out, Complexity.CounterProg.goes_rd_end, Complexity.CounterProg.goes_rd_false, Complexity.UTab.ER, Complexity.UTab.ER.f, Complexity.UTab.ER.zero, Complexity.UTab.Lmax, Complexity.UTab.cellStart, Complexity.UTab.classify, Complexity.UTab.emitC, Complexity.UTab.goes_emit, Complexity.UTab.goes_prefix, Complexity.UTab.goes_rN, Complexity.UTab.goes_rT, Complexity.UTab.goes_stop, Complexity.UTab.lstart, Complexity.UTab.prog, Complexity.UTab.rT, Complexity.UTab.runBound, Complexity.UTab.runBound_le}
\difficulty{6}
Let $M$ be a finite-state Turing machine with Boolean input and let $w$ be a finite
Boolean word. There exist a final register assignment, an input position $p$, an output
word $e$, and a bound $B$ such that $B\le 60(k+1)(L+30)(|w|+1)^3$, where $k$ is the
number of work tapes of $M$ and $L$ is the maximum template length determined by $M$.
Starting at the program's initial label with all eleven registers zero, input position
zero, and any initial output word $o$, the associated counter program reaches a halted
state with the chosen register assignment and input position, appending exactly $e$ to
$o$ in at most $B$ steps.
\end{theorem}

\begin{definition}[Tabulation output word]
\lean{Complexity.UTab.tabEmit}
\label{Complexity.UTab.tabEmit}
\leanok
\uses{Complexity.CounterProg.init, Complexity.CounterProg.run, Complexity.UTab.emitC, Complexity.UTab.prog}
Let $M$ be a finite-tape Turing machine and $w$ a finite Boolean word. Run the program
associated with $M$ on input $w$, from the initial configuration denoted by $rN$, for
$c(M)(|w|+1)^3$ steps, where $c(M)$ is the emission constant; define the tabulation
output to be the Boolean word on the output tape at the end of that run.
\end{definition}

\begin{theorem}[Exact run from a bounded execution]
\lean{Complexity.UTab.run_budget}
\label{Complexity.UTab.run_budget}
\leanok
\uses{Complexity.CounterProg.Goes, Complexity.CounterProg.St, Complexity.CounterProg.init, Complexity.CounterProg.run, Complexity.CounterProg.run_add, Complexity.CounterProg.run_of_halted, Complexity.UTab.ER.f, Complexity.UTab.ER.zero, Complexity.UTab.ER.zero_f, Complexity.UTab.Lb, Complexity.UTab.Lmax, Complexity.UTab.prog}
\difficulty{4}
Let $M$ be a finite Turing machine over the binary alphabet, let $w,e$ be finite binary
lists, let $\rho:\{0,\ldots,10\}\to\mathbb{N}$ be register values, and let
$p,B,C\in\mathbb{N}$. Suppose that for every output list $o$ there is a time $t\le B$ at
which the 11-register program associated with $M$, started at its initial label with all
registers and the input position zero and output $o$, is halted with register values
$\rho$, input position $p$, and output $o$ followed by $e$. If $B\le C$, then after
exactly $C$ steps, starting from that initial configuration with empty output, the
program is halted with register values $\rho$, input position $p$, and output $e$.
\end{theorem}

\begin{theorem}[Polynomial-time computability of the tableau emitter]
\lean{Complexity.UTab.polyTimeComputable_tabEmit}
\label{Complexity.UTab.polyTimeComputable_tabEmit}
\leanok
\uses{Complexity.CounterProg.polyTimeComputable, Complexity.PolyTimeComputable, Complexity.UTab.emitC, Complexity.UTab.goes_any, Complexity.UTab.prog, Complexity.UTab.run_budget, Complexity.UTab.tabEmit}
\difficulty{3}
For every finite Turing machine $M$ over the Boolean alphabet, the function that maps a
Boolean string $w$ to the output produced by the fixed eleven-register program for $M$
after $60(k+1)(L+30)(|w|+1)^3$ steps is polynomial-time computable, where $k$ is the
number of work tapes of $M$ and $L$ is the maximum template length for $M$.
\end{theorem}

\begin{theorem}[Unary tableau output]
\lean{Complexity.UTab.tabEmit_eq}
\label{Complexity.UTab.tabEmit_eq}
\leanok
\uses{BoolCircuit.DAGCircuit.encode, Complexity.UTab.emitC, Complexity.UTab.goes_emit, Complexity.UTab.runBound_le, Complexity.UTab.run_budget, Complexity.UTab.tabEmit, Complexity.cfgTab}
\difficulty{4}
Let $M$ be a finite-state Turing machine with Boolean symbols, let $n,T\in\mathbb{N}$,
and let $z$ be a finite bit string, with $T\geq 1$. The counter program's output on the
input consisting of $n$ ones, a zero, $T$ ones, a zero, and then $z$ is the encoding of
the circuit tableau for $M$, input width $n$, suffix $z$, and time bound $T$.
\end{theorem}
